\documentclass[12pt,a4paper]{report}
\usepackage[utf8]{inputenc}
\usepackage{amsmath}
\usepackage{amsfonts}
\usepackage{amssymb}
\usepackage{amsthm}
\usepackage{hyperref}

\usepackage{multicol}
\usepackage{fancyhdr}
\usepackage[inline]{enumitem}
\usepackage{tikz}
\usepackage{tikz-cd}
\usetikzlibrary{calc}
\usetikzlibrary{shapes.geometric}
\usetikzlibrary{positioning}
\usepackage[margin=0.5in]{geometry}
\usepackage{xcolor}

\hypersetup{
    colorlinks=true,
    linkcolor=blue,
    filecolor=magenta,      
    urlcolor=cyan,
    pdftitle={Tensors},
    pdfpagemode=FullScreen,
    }

%\urlstyle{same}

\newcommand{\CLASSNAME}{Math 5050 -- Special Topics: Tensor Analysis}
\newcommand{\STUDENTNAME}{Paul Carmody}
\newcommand{\ASSIGNMENT}{Chapter 4: Change of Coordinates }
\newcommand{\DUEDATE}{January 2026}
\newcommand{\PROFESSOR}{Professor Berchenko-Kogan}
\newcommand{\SEMESTER}{Spring 2026}
\newcommand{\SCHEDULE}{TBD}
\newcommand{\ROOM}{Remote}

\newcommand{\MMN}{M_{m\times n}}
\newcommand{\FF}{\mathcal{F}}

\pagestyle{fancy}
\fancyhf{}
\chead{ \fancyplain{}{\CLASSNAME} }
%\chead{ \fancyplain{}{\STUDENTNAME} }
\rhead{\thepage}

\fancyfoot{} % clear all footer fields
\fancyfoot[LE,RO]{\thepage}
\fancyfoot[LO,CE]{\ASSIGNMENT}
%\fancyfoot[CO,RE]{To: Dean A. Smith}
\input{../MyMacros}

\newcommand{\RED}[1]{\textcolor{red}{#1}}
\newcommand{\BLUE}[1]{\textcolor{blue}{#1}}

\newcommand{\SUPP}{\operatorname{supp}}
\newcommand{\CLOSURE}{\operatorname{closure of}}
\newcommand{\BD}{\operatorname{bd}}
\newcommand{\HHN}{\mathcal{H}^n}

\begin{document}

\begin{center}
	\Large{\CLASSNAME -- \SEMESTER} \\
	\large{ w/\PROFESSOR}
\end{center}
\begin{center}
	\STUDENTNAME \\
	\ASSIGNMENT -- \DUEDATE\\
\end{center} 

\noindent\textbf{Notes:}

\begin{description}
	\item \textbf{Affine Transformations}
	
	\begin{align*}
		x^i &= A_{i'}^{i}x^{i'}+b^{i} \\
		\PART{x^i}{x^{i'}} &= A_{i'}^{i} \\
		\SPART{x^i}{(x^{i'})} &= 0
	\end{align*}
\end{description}
\HLINE
\noindent\textbf{\Large{Exercises}}

For the following exercises we define
\begin{align*}
	f, A, B, C, A^i &: \R^2 \to \R \implies \text{ e.g., } f(\mu,\nu), A^1(\mu,\nu) \\
	F &: \R^3 \to \R \implies F=F(a^1, a^2, a^3) \\
	f(\mu,\nu) &= F\PAREN{A^1(\mu,\nu), A^2(\mu,\nu), A^3(\mu,\nu)}
\end{align*}
\begin{description}
	\item \textbf{Exercise 36, Pg 41}
	
	Derive the expression for 
	\begin{align*}
		\SPART{f(\mu,\nu)}{\mu}, \frac{\partial^2 f}{\partial\mu\partial\nu}, \AND \SPART{f(\mu,\nu)}{\nu}
	\end{align*}similar to the equations (4.35) and (4.36).
	
	\BLUE{Given the following,\begin{align*}
		\PART{f(\mu,\nu)}{\mu} &= \PART{F}{a^i}\PART{A^i}{\mu} & \quad &\AND & \quad 
		\PART{f(\mu,\nu)}{\nu} &= \PART{F}{a^i}\PART{A^i}{\nu} 
	\end{align*}then
	\begin{align*}
		\PART{}{\mu}\PART{f(\mu,\nu)}{\mu} &= \PART{}{\mu}\PAREN{\PART{F}{a^i}\PART{A^i}{\mu}} \\
		\SPART{f(\mu,\nu)}{\mu} &= \frac{\partial^2 F}{\partial \mu \partial a^i}\PART{A^i}{\mu}+\PART{F}{a^i}\SPART{A^i}{\mu} 
	\end{align*}
	}

%Section 4.6	
	\item \textbf{Exercise 46, Pg 47}  By differentiation
	\begin{align*}
		Z^{i'}(Z(Z'))=Z'
	\end{align*}show that 
	\begin{align*}
		J_i^{i'}J_{j'}^i= \delta^{i'}_{j'}
	\end{align*}
	
	\BLUE{Where
	\begin{align*}
		Z^I(Z'(Z)) \equiv Z^i
	\end{align*}or by identity
	\begin{align*}
		Z^{i'}(Z(Z')) \equiv Z^{i'}
	\end{align*}and
	\begin{align*}
		J_{i'}^i &= \PART{Z^i(Z)}{Z^{i'}} & \quad
		J_i^{i'} &= \PART{Z^{i'}(Z)}{Z^i} \\
			&= \PART{Z^i}{Z}\PART{Z}{Z^{i'}} & \quad &= \PART{Z^{i'}}{Z}\PART{Z}{Z^i}
	\end{align*}which follows that
	\begin{align*}
		J_i^{i'}J_{j'}^i &= \PART{Z^{i'}(Z)}{Z^i}\PAREN{\PART{Z^i(Z)}{Z^{i'}}} \\
		&= \PART{Z^{i'}}{Z}\PART{Z}{Z^i}\PAREN{\PART{Z^i}{Z}\PART{Z}{Z^{j'}}} \\
		&= \PART{Z^{i'}}{Z^{j'}} \\
		&= \delta_{j'}^{i'}
	\end{align*}
	}
	
	\item \textbf{Exercise 47, Pg 47} Derive (4.76) from (4.74) by multiiplying both sides by $J_{j'}^j$.
	
	\BLUE{
	\begin{align*}
		J_{j'}^j &= \PART{Z^{j}(Z)}{Z^{j'}} = \PART{Z^j}{Z}\PART{Z}{Z^{j'}}\\
		J_i^{i'}J_{i'}^j&= \delta^{i}_{j} & (4.74) \\
		J_i^{i'}J_{i'}^jJ_{j'}^j &= \delta^{i}_{j}J_{j'}^j \\
		\PAREN{\PART{Z^{i'}}{Z}\PART{Z}{Z^i}\PAREN{\PART{Z^j}{Z}\PART{Z}{Z^{i'}}}}\PART{Z^j}{Z}\PART{Z}{Z^{j'}} &= \delta^{i'}_{j'}J_{j'}^j \\
		\PAREN{\PART{Z^{i'}}{Z^i}\PAREN{\PART{Z^j}{Z^{i'}}}}\PART{Z^j}{Z}\PART{Z}{Z^{j'}} &= \delta^{i'}_{j'}J_{j'}^j \\
		\PART{Z^j}{Z^i}\PART{Z^j}{Z^{j'}} &= \delta^{i'}_{j'}J_{j'}^j
	\end{align*}
	}
	
	\item \textbf{Exercise 48, Pg 47} Introduce the symbols $J_{i'j'}^i$ and $J_{ij}^{i'}$ for the second order partial derivatives
	\begin{align*}
		J_{i'j'}^i &= \frac{\partial^2 Z^i(Z')}{\partial Z^{i'}\partial Z^{j'}}\\
		J_{ij}^{i'} &= \frac{\partial^2 Z^{i'}(Z)}{\partial Z^{i}\partial Z^{j}}.
	\end{align*}Show that
	\begin{align*}
		J_{i'j'}^iJ_j^{i'}J_k^{j'}+J_{i'}^iJ_{jk}^{i'}=0.
\end{align*}	 How many identities does this tensor relationship represent?
	
	\BLUE{Let's begin with 
	\begin{align*}
		J_i^{i'}J_{j'}^i &=\delta_{j'}^{i'} \\
		\frac{\partial}{\partial Z^k}\delta_{j'}^{i'} &=\frac{\partial}{\partial Z^k}\PAREN{J_i^{i'}J_{j'}^i} \\
		0 &= \frac{\partial}{\partial Z^k}\PAREN{J_i^{i'}}J_{j'}^i + J_i^{i'}\frac{\partial}{\partial Z^k}\PAREN{J_{j'}^i}\\
		&= J_{ik}^{i'}J_{j'}^i+J_i^{i'}J_{kj'}^i
	\end{align*}
	}
	
	\newcommand{\JJJ}[2]{J^{#1}_{#2}}
	\newcommand{\JJJJ}[2]{\frac{\partial Z^{#1}(Z)}{\partial Z^{#2}}}
	\item \textbf{Exercise 53, Pg 47} Consider the three changes of variable: $Z^i$, $Z^{i'}$, $Z^{i''}$ and back to $Z^i$.  Show that
	\begin{align*}
		\JJJ{i}{i'}\JJJ{i'}{i''}\JJJ{i''}{j}	&= \delta^i_j
	\end{align*}
	
	\BLUE{\begin{align*}
		\JJJ{i}{i'}\JJJ{i'}{i''}\JJJ{i''}{j}	&= \JJJJ{i}{i'}\JJJJ{i'}{i''}\JJJJ{i''}{j} \\
		&= \JJJJ{i}{}\PART{Z}{Z^{i'}}\JJJJ{i'}{}\PART{Z}{Z^{i''}}\JJJJ{i''}{}\PART{Z}{Z^j} \\
		&= \PART{Z^i(Z)}{Z^{i'}}\PART{Z^{i'}(Z)}{Z^{i''}}\PART{Z^{i''}(Z)}{Z^{j}} \\
		&= \PART{Z^i(Z)}{Z^{j}} \\
		&=  \delta^i_j
	\end{align*}
	}

#Section 4.9	
	\item \textbf{Exercise 54, Pg 50} Simplify $A_{ij}\delta^i_j$.
	
	\BLUE{\begin{align*}
		A_{ij}\delta^i_j &= \sum_{i=1}^n A_{ij}\delta^i_j = A_{jj}
	\end{align*}
	}
	
	\item \textbf{Exercise 55, Pg 50} Simplify $\delta^i_{j}\delta^j_k$.
	
	\BLUE{\begin{align*}
		\delta^i_{j}\delta^j_k &= \sum_{j=1}^n \delta^i_{j}\delta^j_k = \delta^i_k
	\end{align*}
	}
	
	\item \textbf{Exercise 56, Pg 50} Simplify $\delta^i_{j}\delta^j_i$.
	
	\BLUE{\begin{align*}
		\delta^i_{j}\delta^j_i &= \sum_{i=1}^n\sum_{j=1}^n \delta^i_{j}\delta^j_i = \delta^i_i = \sum_{i=1}^n 1 = n
	\end{align*}This is the trace of the identity matrix/tensor.
	}
	
	\item \textbf{Exercise 57, Pg 50} Simplify $\delta^i_{i}\delta^j_j$.
	
	\BLUE{\begin{align*}
		\delta^i_{i}\delta^j_j &= \sum_{i=1}^n\sum_{j=1}^n \delta^i_{i}\delta^j_j = \sum_{i=1}^n\delta^i_{i}\sum_{j=1}^n \delta^j_j = n\cdot n
	\end{align*}
	}
	
	\item \textbf{Exercise 58, Pg 50} Simplify $\delta^i_{j}S_iS^j$.
	
	\BLUE{\begin{align*}
		\delta^i_{j}S_iS^j &= \sum_{i=1}^n\sum_{j=1}^n \delta^i_{j}S_iS^j = \sum_{i=1}^n S_iS^i = S_iS^i = |S|
	\end{align*}
	}
	
	\item \textbf{Exercise 58, Pg 50} Show that the object $\delta^r_i\delta^s_j-\delta^s_i\delta^r_j$ is skew-symmetric in $i$ and $j$ as well as $r$ and $s$.
	
	\BLUE{Let $A^{rs}_{ij}=\delta^r_i\delta^s_j-\delta^s_i\delta^r_j$.  Then $A^{rs}_{ij}$ is skew-symmetric in $i$ and $j$ if $A^{rs}_{ij}=-A^{rs}_{ji}$.  Let $B_{ij} = A^{rs}_{ij}$ then
	\begin{align*}
		B_{ij}+B_{ji} &= \PAREN{\delta^r_i\delta^s_j-\delta^s_i\delta^r_j}-\PAREN{\delta^r_j\delta^s_i-\delta^s_j\delta^r_i} ,\, \forall r,s \in \{1,\dots,n\} \\
		&= \PAREN{\delta^r_i\delta^s_j-\delta^s_i\delta^r_j-\delta^r_j\delta^s_i+\delta^s_j\delta^r_i}\\
		&= \PAREN{\delta^r_i\delta^s_j-\delta^r_j\delta^s_i} +\PAREN{\delta^s_i\delta^r_j-\delta^s_j\delta^r_i} \\
		&= 0
\end{align*}	Therefore, $A^{rs}_{ij}= -A^{rs}_{ji}$.  Similarly, by setting $B^{rs}=A^{rs}_{ij}$.
	}
\end{description}
\end{document}
